% F085｜真实数学图；代表函数不构成抽象解的唯一性证明。
\begin{center}
\begin{tikzpicture}
\begin{axis}[width=12.4cm,height=5.1cm,axis lines=middle,
 xmin=-0.3,xmax=6.7,ymin=-2,ymax=41,
 xtick={0,1,2,3,6},xticklabels=\empty,ytick={0,4,9,36},xlabel={$t$},ylabel={$f(t)$},
 title={幂函数曲线 $f(t)=t^2$（这里取 $t>0$）},title style={font=\small,yshift=3pt},
 tick label style={font=\small},label style={font=\small},
 ylabel style={at={(axis description cs:0,1)},anchor=south east,xshift=-18pt},
 clip=false,grid=major,grid style={black!8},enlargelimits=false,
 every axis plot/.append style={thick},samples=161]
\addplot[CExplanation,very thick,domain=0:6.2] {x^2};
\draw[dashed,black!35] (axis cs:2,0) |- (axis cs:0,4);
\draw[dashed,black!35] (axis cs:3,0) |- (axis cs:0,9);
\draw[dashed,black!35] (axis cs:6,0) |- (axis cs:0,36);
\addplot[only marks,mark=*,mark size=2pt,orange!80!black] coordinates {(2,4) (3,9) (6,36)};

\node[font=\small,anchor=north,yshift=-3pt] at (axis cs:1,-2) {$1$};
\node[font=\small,anchor=north,yshift=-3pt] at (axis cs:2,-2) {$2$};
\node[font=\small,anchor=north,yshift=-3pt] at (axis cs:3,-2) {$3$};
\node[font=\small,anchor=north,yshift=-3pt] at (axis cs:6,-2) {$6$};

\end{axis}
\end{tikzpicture}
\par\small 横坐标相乘 $2\cdot3=6$，高度也相乘 $4\cdot9=36$；$(0,0)$ 仅是边界参照。
\end{center}
